% Generated draft: requires mathematical and source review.
\documentclass[11pt]{article}
\usepackage[margin=25mm]{geometry}
\usepackage{fontspec}
\setmainfont{Latin Modern Roman}
\usepackage{amsmath,amssymb,mathtools,unicode-math}
\setmathfont{Latin Modern Math}
\usepackage{enumitem,needspace,xurl,hyperref}
\setlength{\parindent}{0pt}
\setlength{\parskip}{5pt}
\setlength{\emergencystretch}{3em}
\newcommand{\N}{\mathbb{N}}
\newcommand{\Z}{\mathbb{Z}}
\newcommand{\Q}{\mathbb{Q}}
\newcommand{\R}{\mathbb{R}}
\newcommand{\C}{\mathbb{C}}
\newcommand{\F}{\mathbb{F}}
\newcommand{\A}{\mathbb{A}}
\newcommand{\PP}{\mathbb{P}}
\newcommand{\E}{\mathbb{E}}
\newcommand{\eps}{\varepsilon}
\newcommand{\dd}{\,\mathrm{d}}
\begin{document}
\section*{UnsolvedMath ID: 701619}
\textbf{Source problem code:} catalogue entry
\subsection{Definitions and mathematical statement}
For a topological group \(G\), a compact open subgroup is a subgroup \(U\leq G\) that is compact and open in the topology of \(G\). An element \(g\in G\) normalizes \(U\) if \(gUg^{-1}=U\); equivalently, \(g\in N_G(U)\), where
\[
N_G(U)=\{x\in G:xUx^{-1}=U\}.
\]
The source asks whether there exists a group \(G\) in its class \(\mathscr{S}\) such that
\[
\forall g\in G\;\exists U_g\leq G:\quad
U_g\text{ is compact and open in }G
\quad\text{and}\quad gU_gg^{-1}=U_g.
\]
The subgroup \(U_g\) may depend on \(g\); a single compact open subgroup normalized by every element is not required.

\textbf{Unresolved qualification.} The supplied record does not define \(\mathscr{S}\) or explicitly specify the topological hypotheses on its members. The displayed existence question therefore retains this undefined source class as a parameter. Its intended domain must be established before a complete mathematical formulation can be certified.
\subsection{Short English statement}
Does the source class \(\mathscr{S}\) contain a group in which each element normalizes some compact open subgroup, possibly depending on that element? The supplied record leaves \(\mathscr{S}\) undefined.
\subsection{Sources}
\begin{itemize}
\item Supplied local registry record: permanent UnsolvedMath ID \(701619\), source code \texttt{catalogue entry}, title \emph{Kourovka Notebook Problem 19.71}. Its status is supplied as \texttt{open}; this label has not been independently verified.
\item Supplied source attribution: \emph{The Kourovka Notebook (Unsolved Problems in Group Theory)}, 19th issue (2018), problem 19.71; also identified with the 21st issue edition (2026), \texttt{arXiv:1401.0300v48}. The collection description names E. I. Khukhro and V. D. Mazurov as editors. These bibliographic details are supplied and have not been independently checked.
\item The background and proposer field attribute the problem to C. Reid, with proposed year 2018. The statement itself begins with the attribution ``(G. Willis).'' Both attributions are retained; their respective roles are not explained in the record.
\end{itemize}
\end{document}
